\documentclass[11pt]{article}
\usepackage[T1]{fontenc}
\usepackage{lmodern}
\usepackage[margin=1in]{geometry}
\usepackage{amsmath,amssymb,amsthm,booktabs}
\usepackage{tabularx}
\usepackage[hidelinks]{hyperref}
\hypersetup{pdftitle={OEIS A209723: Barker's conjectures on distinct clockwise edge differences},
pdfauthor={Haobo Ma and Wenlin Zhang},pdfsubject={CC BY 4.0; exact enumeration of ternary arrays}}
\newtheorem{theorem}{Theorem}
\newtheorem{lemma}{Lemma}
\newtheorem{corollary}{Corollary}
\title{OEIS A209723: Barker's conjectures on\texorpdfstring{\\}{ }distinct clockwise edge differences}
\author{Haobo Ma\\{\small ChronoAI Pte Ltd, Singapore}\and
Wenlin Zhang\\{\small National University of Singapore}\\{\small ChronoAI Pte Ltd, Singapore}}
\date{November 2, 2026}
\begin{document}
\maketitle
\begin{abstract}
We prove the ``Conjectures'' of Colin Barker in the FORMULA section of
OEIS A209723, dated July 12, 2018: the parity formulas and ordinary generating
function for ternary arrays with distinct clockwise edge differences
\cite{a209723}. We also prove Hardin's empirical recurrence. A decomposition
of the row-compatibility graph into two three-vertex stars and isolated edges
gives a formula for every rectangular size. Its corollaries include the
companion formulas in A209721--A209726, Beck's first A209721 conjecture,
Peterson's A209721 recurrence and triangle interpretation, and the column
recurrences and diagonal identification in A209727.
\end{abstract}

\section{The enumeration problem}
OEIS A209723 asks for the normalized count of five-column ternary arrays
with distinct signed clockwise edge differences; Barker's July 12, 2018
FORMULA conjectures give its parity formulas and generating function
\cite{a209723}. We prove them by counting all rectangular sizes at once.
Let $n,k\geq1$. Write $N(n,k)$ for the number of labelled
$(n+1)\times(k+1)$ arrays with entries in $\{0,1,2\}$ such that every consecutive
subblock $\left[\begin{smallmatrix}a&b\\c&d\end{smallmatrix}\right]$ has four
pairwise distinct signed integer differences
\[
 b-a,\qquad d-b,\qquad c-d,\qquad a-c.
\]
Set $T(n,k)=N(n,k)/4$, using Hardin's normalization in
A209727 \cite{a209727}. No quotient by rotations or relabellings is taken.
In particular, A209723 is $a(n)=T(n,4)$, beginning $6,7,8,10,12,16$.

For $m\geq0$, define
\begin{equation}\label{eq:b}
 b(2m+1)=2^{m+1},\qquad b(2m+2)=3\cdot2^m.
\end{equation}
Thus $b(n)$ is A029744 at index $n+1$ \cite{a029744}.
\begin{theorem}\label{thm:main}
For all $n,k\geq1$,
\begin{equation}\label{eq:main}
 T(n,k)=b(n)+b(k)-2.
\end{equation}
Consequently,
\begin{equation}\label{eq:target}
 a(2m+1)=2^{m+1}+4,\qquad a(2m+2)=3\cdot2^m+4\quad(m\geq0),
\end{equation}
and $\sum_{n\geq1}a(n)x^n=x(6+x-11x^2)/((1-x)(1-2x^2))$.
\end{theorem}

The source calls the recurrence ``Empirical'' and attributes the generating
function and parity formulas to Barker as conjectures \cite{a209723}.
For the four-column case A209722, Andrew's 2024 Math Stack Exchange answer
already gives a tile-based proof outline \cite{andrew}.
Here the graph decomposition yields a complete proof for all widths at once.
The underlying correspondence between arrays and walks is the usual
transfer-matrix method; see Stanley \cite[Section 4.7]{stanley}.
We give the required counting argument directly.

\section{Two stars and isolated edges}
For a fixed $k$, let $H_k$ have the words in $\{0,1,2\}^{k+1}$ as vertices.
Two words $s,t$ are adjacent precisely when using them as consecutive rows
gives valid subblocks in all $k$ positions. Reversing the two rows permutes
and negates the differences, so this relation is symmetric. Loops are
impossible: a repeated row would give two zero vertical differences.
An array counted by $N(n,k)$ is uniquely a walk of length $n$ in $H_k$,
with a specified starting vertex. Walks may revisit vertices and edges.

For a subblock, the differences lie in $\{-2,-1,0,1,2\}$ and sum to zero.
If four are distinct, their omitted fifth value must therefore be zero.
Both a $+2$ and a $-2$ edge occur. These cannot be opposite edges of the
square: if they were, the two remaining differences would coincide.
Thus they are adjacent and share an extreme entry. The remaining vertex
must have value $1$ to supply differences $+1$ and $-1$.
This gives exactly the following eight tiles, listed as $(a,b,c,d)$:
\begin{equation}\label{eq:tiles}
\begin{split}
 &(0,1,2,0),\ (0,2,1,0),\ (0,2,2,1),\ (1,0,0,2),\\
 &(1,2,2,0),\ (2,0,0,1),\ (2,0,1,2),\ (2,1,0,2).
\end{split}
\end{equation}
Each listed tile satisfies the condition, which also verifies sufficiency.

Read a two-row array by columns. The nonisolated vertices of $H_1$ are the
ordered pairs $01,02,10,12,20,21$, with adjacency
\begin{equation}\label{eq:stars}
 02\longleftrightarrow\{10,21\},\qquad
 20\longleftrightarrow\{01,12\}.
\end{equation}
This is the disjoint union of two copies of the star $K_{1,2}$.
In particular, the neighbors of any column state have distinct first entries.

\begin{lemma}\label{lem:structure}
For every $k\geq1$, the nonisolated part of $H_k$ consists of two copies
of $K_{1,2}$ and $2b(k)-4$ copies of $K_2$.
\end{lemma}
\begin{proof}
For distinct symbols $u,v$, let $R_{uv}$ be the length-$(k+1)$ word
$uvuvu\cdots$. The two words $R_{02}$ and $R_{20}$ each have exactly two
neighbors. For $R_{02}$, if the first column is $01$, the next must be $20$;
if it is $02$, the next first entry $2$ forces the next column to be $21$.
Continuing with \eqref{eq:stars} determines the whole lower row.
The case $R_{20}$ follows by interchanging $0$ and $2$.
The resulting two stars are
\[
 R_{02}\longleftrightarrow\{R_{10},R_{21}\},\qquad
 R_{20}\longleftrightarrow\{R_{01},R_{12}\}.
\]
The six words are distinct since $k+1\geq2$.

To show that there are no other branching vertices, take a row $s$ with
at least one neighbor. Its adjacent entries are unequal by
\eqref{eq:tiles}. If $s$ contains no $1$, it is $R_{02}$ or $R_{20}$.
Otherwise some adjacent pair of $s$ involves $1$, since the row has at
least two entries. Inspection of \eqref{eq:tiles} shows that each such
top pair has exactly one permitted lower pair. Starting with that pair,
extend to the right and left. At each step the top entry of the next
column is prescribed, and \eqref{eq:stars} supplies at most one neighboring
column with that first entry. The extension is therefore unique whenever
it exists. All rows other than the two centers have degree at most one.
The four displayed leaves have degree one. Removing these two stars
therefore leaves a graph of maximum degree one, whose nonisolated
components are copies of $K_2$.

It remains to count these edges. Reading the same two-row array either
by rows or by columns shows that the number of directed edges of $H_k$
equals the number of walks of length $k$ in $H_1$.
In a single $K_{1,2}$ the total number of walks of length $2m$ is
$3\cdot2^m$: each of the three starting vertices has $2^m$ such walks.
For length $2m+1$, the center contributes $2^{m+1}$ and each leaf
contributes $2^m$, for a total $4\cdot2^m$.
Thus the two stars together have $4b(k)$ walks of length $k\geq1$.
The two stars in $H_k$ account for eight directed edges. Each remaining
$K_2$ accounts for two, giving $(4b(k)-8)/2=2b(k)-4$ copies.
\end{proof}

\begin{proof}[Proof of Theorem \ref{thm:main}]
For positive walk length $n$, isolated vertices contribute nothing.
By the star calculation, the two stars of $H_k$ contribute $4b(n)$
walks. Each of the $2b(k)-4$ copies of $K_2$ contributes exactly two:
there is one walk from either starting vertex. Therefore
\[
 N(n,k)=4b(n)+2(2b(k)-4)=4\bigl(b(n)+b(k)-2\bigr).
\]
For $k=4$, $b(4)=6$, proving \eqref{eq:target}.
The generating function follows from
\begin{equation}\label{eq:gf}
 B(x):=\sum_{n\geq1}b(n)x^n=\frac{2x+3x^2}{1-2x^2},\qquad
 \sum_{n\geq1}a(n)x^n=B(x)+\frac{4x}{1-x}.
\end{equation}
These identities hold as formal power series, since the denominators
have constant term one.
\end{proof}

\section{Companion formulas and source coverage}
Put $C_k=b(k)-2$ and $A_k(x)=\sum_{n\geq1}T(n,k)x^n$.
The same calculation gives, for every $k\geq1$,
\begin{align}
 T(2m+1,k)&=2^{m+1}+C_k,&
 T(2m+2,k)&=3\cdot2^m+C_k,\label{eq:companions}\\
 A_k(x)&=\frac{x\bigl((C_k+2)+x-(2C_k+3)x^2\bigr)}
 {(1-x)(1-2x^2)}.\label{eq:generalgf}
\end{align}
Also $T(n+2,k)=2T(n,k)-C_k$ for $n\geq1$.
Subtracting this identity at two consecutive indices proves
\begin{equation}\label{eq:recurrence}
 T(n,k)=T(n-1,k)+2T(n-2,k)-2T(n-3,k)\qquad(n\geq4).
\end{equation}
For $k=1$, $C_1=0$, so the stronger recurrence is
$T(n,1)=2T(n-2,1)$ for $n\geq3$.
These prove every column recurrence listed in A209727, not just those
listed for $k\leq7$ \cite{a209727}.
For each fixed $k$, the same parity formulas also imply
$\lim_{n\to\infty}T(n,k)^{1/n}=\sqrt{2}$: on either parity the
exponential factor is a constant multiple of $2^{n/2}$, and the
additive constant $C_k$ does not change its $n$th-root limit.

\begin{corollary}\label{cor:beck-peterson}
For $a(n)=T(n,2)$, the sequence A209721,
\[
 a(n)=\mathrm{A029744}(n+1)+1\quad(n\geq1),\qquad
 a(n)=2a(n-2)-1\quad(n\geq3).
\]
For $n\geq3$, $a(n)$ also counts integer triangles having two specified
side lengths $a(n-1)$ and $a(n-2)$.
\end{corollary}
\begin{proof}
Since $b(2)=3$, Theorem~\ref{thm:main} gives $a(n)=b(n)+1$.
The identification $b(n)=\mathrm{A029744}(n+1)$ proves Beck's first
March 26, 2022 conjecture. Equation~\eqref{eq:b} gives
$b(n)=2b(n-2)$ for $n\geq3$, so
$a(n)=2(a(n-2)-1)+1=2a(n-2)-1$, proving Peterson's April 26, 2020
empirical recurrence \cite{a209721}.
The parity formulas show that $a(n)$ is strictly increasing.
Writing $A=a(n-1)$ and $B=a(n-2)$, the third side is any integer from
$A-B+1$ to $A+B-1$, inclusive. There are $2B-1=a(n)$ choices.
\end{proof}

\begin{corollary}\label{cor:diagonal}
The diagonal of A209727 satisfies
\[
 T(n,n)=2b(n)-2=\mathrm{A027383}(n)\qquad(n\geq1).
\]
\end{corollary}
\begin{proof}
Set $k=n$ in Theorem~\ref{thm:main}. For $n=2m\geq2$, the result is
$3\cdot2^m-2$; for $n=2m+1\geq1$, it is $2^{m+2}-2$.
These are the defining parity formulas of A027383 \cite{a027383}, proving
the ``Diagonal is A027383'' cross-reference in A209727 \cite{a209727}.
The OEIS offset of A027383 is zero, but the array diagonal begins at
$n=1$, with $2,4,6,10,14,\ldots$; there is no index shift or assertion
about $T(0,0)$.
\end{proof}

The following coverage map distinguishes the source statements.
For A209721 the generating function is labelled empirical; Barker's
formulas in A209722--A209726 are labelled conjectures.
\begin{center}
\small
\begin{tabularx}{\textwidth}{@{}p{0.19\textwidth}X p{0.26\textwidth}@{}}
\toprule
Source & Statement & Coverage\\
\midrule
A209721 \cite{a209721} & Empirical order-three recurrence and Barker's
March 23, 2012 generating function ($k=2$, $C_2=1$)
& Proved by \eqref{eq:generalgf}--\eqref{eq:recurrence}\\
A209721 & Beck's first conjecture; Peterson's recurrence and triangle count
& Proved by Corollary~\ref{cor:beck-peterson}\\
A209721 & Beck's second conjecture: zero positions in
$\mathrm{A309019}(n+2)-\mathrm{A002487}(n+2)$ & Not covered\\
A209722 \cite{a209722} & Barker's July 12, 2018 formulas and empirical
recurrence ($k=3$, $C_3=2$) & Proved here; Andrew's earlier outline
\cite{andrew} cited\\
A209723 \cite{a209723} & Barker's July 12, 2018 formulas and empirical
recurrence ($k=4$, $C_4=4$) & Proved by Theorem~\ref{thm:main} and
\eqref{eq:recurrence}\\
A209724 \cite{a209724} & Barker's March 7, 2018 formulas and empirical
recurrence ($k=5$, $C_5=6$) & Proved by \eqref{eq:companions}--\eqref{eq:recurrence}\\
A209724 & Reverter i Bigas's October 11, 2020 smallest-integer/divisor
conjecture & Not covered\\
A209725 \cite{a209725} & Barker's July 12, 2018 formulas and empirical
recurrence ($k=6$, $C_6=10$) & Proved by \eqref{eq:companions}--\eqref{eq:recurrence}\\
A209726 \cite{a209726} & Barker's July 12, 2018 formulas and empirical
recurrence ($k=7$, $C_7=14$) & Proved by \eqref{eq:companions}--\eqref{eq:recurrence}\\
A209727 \cite{a209727} & Column recurrences; column 1 is A029744$(n+1)$
& Proved above for every $k\geq1$\\
A209727 & Diagonal is A027383 & Proved by Corollary~\ref{cor:diagonal}\\
\bottomrule
\end{tabularx}
\end{center}
The formulas here start at $n=1$; assigning a zero constant term to a
generating function does not extend the array enumeration to $n=0$.
The accompanying Python script independently enumerates all $81$ tiles,
compares row-pair enumeration with column-walk enumeration for widths
$2$--$6$, and checks the component decomposition through width $12$.
It checks all $210$ tabulated A209723 values and all $1512$ tabulated
A209727 values. These finite checks support, but do not replace, the
all-dimension proof above.

\paragraph{Lean statement and proof scope.}
The declaration \texttt{result} in the
\href{https://github.com/the-omega-institute/trureturing/blob/cdafeab0d44da2b6840cfc802535d2f94a146fad/D5/S3/Combinatorics/Graph/Pressure/HardinBarkerA209723EdgeDifferences.lean}{pinned five-column Lean file}
formalizes the case $k=4$: the two parity identities,
the recurrence for $n\geq4$, existence of the zero-constant source series,
and its denominator identity. Its array entries are in $\mathrm{Fin}\,3$,
its differences are signed integers, and its count is divided by four,
as in our definition of $T(n,4)$.
The all-width theorem and Corollaries~\ref{cor:beck-peterson}
and~\ref{cor:diagonal} have the written proofs above; the Python checks
cover finite instances of those statements.

\paragraph{License and reproducibility.}
Copyright 2026 Haobo Ma and Wenlin Zhang.
This note is licensed under Creative Commons Attribution 4.0 International
(CC BY 4.0). The accompanying source package contains this LaTeX document,
the verification script, and build instructions. No external computation
is needed to check the proof.

\begin{thebibliography}{10}
\bibitem{a209723} R. H. Hardin, OEIS A209723: five-column ternary arrays with
distinct clockwise edge differences, March 12, 2012; FORMULA contributions
by C. Barker, July 12, 2018. \url{https://oeis.org/A209723}.
\bibitem{a209727} R. H. Hardin, OEIS A209727: rectangular ternary arrays with
distinct clockwise edge differences, March 12, 2012.
\url{https://oeis.org/A209727}.
\bibitem{a029744} N. J. A. Sloane, OEIS A029744: Numbers of the form
$2^k$ or $3\cdot2^k$. \url{https://oeis.org/A029744}.
\bibitem{andrew} Andrew, answer to ``Could someone provide an analytical
proof of empirically found formula?'', Math Stack Exchange, November 8, 2024
(00:07:06 UTC),
answer 4995628. \url{https://math.stackexchange.com/a/4995628}.
\bibitem{stanley} R. P. Stanley, \emph{Enumerative Combinatorics, Volume 1},
second edition, Cambridge University Press, 2012, Section 4.7.
\url{https://doi.org/10.1017/CBO9781139058520}.
\bibitem{a209721} R. H. Hardin, OEIS A209721: three-column arrays,
March 12, 2012; formula by C. Barker, March 23, 2012.
\url{https://oeis.org/A209721}.
\bibitem{a209722} R. H. Hardin, OEIS A209722: four-column arrays,
March 12, 2012; conjectures by C. Barker, July 12, 2018.
\url{https://oeis.org/A209722}.
\bibitem{a209724} R. H. Hardin, OEIS A209724: six-column arrays,
March 12, 2012; conjectures by C. Barker, March 7, 2018.
\url{https://oeis.org/A209724}.
\bibitem{a209725} R. H. Hardin, OEIS A209725: seven-column arrays,
March 12, 2012; conjectures by C. Barker, July 12, 2018.
\url{https://oeis.org/A209725}.
\bibitem{a209726} R. H. Hardin, OEIS A209726: eight-column arrays,
March 12, 2012; conjectures by C. Barker, July 12, 2018.
\url{https://oeis.org/A209726}.
\bibitem{a027383} R. K. Guy, OEIS A027383:
$a(2m)=3\cdot2^m-2$, $a(2m+1)=2^{m+2}-2$.
\url{https://oeis.org/A027383}.
\end{thebibliography}
\end{document}
